\DocumentMetadata{
  pdfversion=2.0,
  pdfstandard=ua-2,
  lang=en-US,
  tagging=on,
  tagging-setup={math/setup=mathml-SE,tabsorder=structure}
}
\documentclass[11pt,letterpaper]{article}
\usepackage[margin=0.68in,includefoot]{geometry}
\usepackage{newcomputermodern}
\usepackage{amsmath,amscd,mathtools}
\usepackage{tikz,adjustbox}
\providecommand{\square}{\mdlgwhtsquare}
\usepackage[hidelinks]{hyperref}
\hypersetup{
  pdftitle={Analysis PhD Exam, September 2002},
  pdfauthor={University of Florida Department of Mathematics},
  pdfsubject={Graduate mathematics examination},
  pdfdisplaydoctitle=true,
  bookmarksopen=true
}
\setcounter{secnumdepth}{0}
\setlength{\parindent}{0pt}
\setlength{\parskip}{0.28em}
\setlength{\emergencystretch}{3em}
\setlength{\leftmargini}{2.15em}
\setlength{\leftmarginii}{2.65em}
\setlength{\leftmarginiii}{3em}
\renewcommand{\labelenumi}{\textbf{\theenumi.}}
\pagestyle{empty}
\raggedbottom
\allowdisplaybreaks
\newenvironment{examproblems}[1]{%
  \begin{enumerate}%
  \setcounter{enumi}{#1}%
  \setlength{\topsep}{0.35em}%
  \setlength{\partopsep}{0pt}%
  \setlength{\itemsep}{0.48em}%
  \setlength{\parsep}{0.18em}%
}{\end{enumerate}}
\newenvironment{examsubparts}{%
  \begin{enumerate}%
  \setlength{\topsep}{0.18em}%
  \setlength{\partopsep}{0pt}%
  \setlength{\itemsep}{0.16em}%
  \setlength{\parsep}{0.1em}%
}{\end{enumerate}}
\newenvironment{examinstructions}{%
  \par\begingroup\itshape
}{%
  \par\endgroup\vspace{0.18em}
}
\makeatletter
\renewcommand\section{\@startsection{section}{1}{0pt}%
  {-1.25ex plus -.25ex minus -.1ex}%
  {0.55ex plus .1ex}%
  {\normalfont\large\bfseries}}
\renewcommand{\@maketitle}{%
  \newpage\null\vskip -1em%
  {\footnotesize\bfseries
    \href{https://www.ufl.edu/}{University of Florida}, \href{https://math.ufl.edu/}{Department of Mathematics}\par}%
  \vskip -0.5em%
  \tagstructbegin{tag=Title}%
    {\LARGE\bfseries\href{https://gma.math.ufl.edu/exams/analysis-phd/}{Analysis PhD Exam}, September 2002\par}%
  \tagstructend%
  \vskip 0.25em%
  \tagmcbegin{artifact}%
    \hrule height 1pt%
  \tagmcend%
  \vskip 1em%
}
\makeatother
\title{Analysis PhD Exam, September 2002}
\author{}
\date{}
\begin{document}
\maketitle
\thispagestyle{empty}
\section{I}
\begin{examinstructions}
State the following theorems.
\end{examinstructions}
\begin{examproblems}{0}
\item
The Egorov theorem.
\item
The Lebesgue theorem.
\item
The dual of \(L^p(\mu)\), \(1\leq p<\infty\).
\item
The Fubini theorem for \(\mu\times\nu\)-integrable functions.
\item
The Riesz representation theorem (of linear continuous functionals on the space \(C(K)\) of continuous functions on a compact space~\(K\)).
\end{examproblems}
\section{II}
\begin{examinstructions}
State and prove one of the following 2 theorems.
\end{examinstructions}
\begin{examproblems}{0}
\item
The representation of the dual of \(L^1(\mu)\).
\item
The Fubini theorem for \(\mathcal S\times\mathcal T\)-measurable functions.
\end{examproblems}
\section{III}
\begin{examinstructions}
Solve the following problems.
\end{examinstructions}
\begin{examinstructions}
Let \((X,\Sigma,\mu)\) be a measure space and~\(E\) a Banach space.
\end{examinstructions}
\begin{examproblems}{0}
\item
Assume~\(\mu\) is finite. Let \(\mathcal F\subset\Sigma\) be a sub-\(\sigma\)-algebra. Prove the existence of the conditional expectation \(E(f\mid\mathcal F)\) for functions \(f\in L_E^1(\mu)\).
\item
Let \(\mathcal R\subset\Sigma\) be an algebra generating~\(\Sigma\) and \(f\in L_E^1(\mu)\). Prove that there is a sequence \((f_n)\) of \(\mathcal R\)-step functions converging to~\(f\), \(\mu\)-a.e., and in the mean.
\item
Let \(\mathcal R\subset\Sigma\) be an algebra generating~\(\Sigma\) and \(f\in L_E^1(\mu)\). Prove that if \(\int_A f\,d\mu=0\) for every \(A\in\mathcal R\), then \(f=0\), \(\mu\)-a.e.
\item
Assume~\(\mu\) is finite and let \(f_n\in L_E^1(\mu)\), \(n=1,2,\ldots\). Prove that if \(f_n\to f\) uniformly, then \(f\in L_E^1(\mu)\) and \(f_n\to f\) in the mean.
\item
Let \((X,\mathcal S)\), \((Y,\mathcal T)\) be measurable spaces, and let \(\mu:\mathcal S\to\mathbb R\) and \(\nu:\mathcal T\to\mathbb R\) be~\(\sigma\)-additive measures. Prove the existence of the product measure \(\mu\times\nu\) and the equality \(\lvert\mu\times\nu\rvert=\lvert\mu\rvert\times\lvert\nu\rvert\).
\end{examproblems}
\begin{examinstructions}
Note. Write complete computations and explanations. Do not use arrows or other unnecessary signs. Use words for explanations. Write complete sentences. Write complete statements of theorems (including the framework).
\end{examinstructions}
\end{document}
